\chapter{Getting Started}

\section{Exercise 2.2-1: Express a function using $\Theta$-notation}

For
\begin{equation*}
    f(n) = \frac{n^3}{1000} - 100n^2 - 100n + 3,
\end{equation*}
the cubic term dominates the lower-order terms as $n$ grows. In particular,
\begin{equation*}
    \lim_{n \to \infty} \frac{f(n)}{n^3} = \frac{1}{1000} > 0.
\end{equation*}
Therefore, $f(n) = \Theta(n^3)$. The coefficient $1/1000$ is ignored when
expressing asymptotic growth; it is not changed to $1/10$ in the original
function.

\section{Exercise 2.2-2: Selection sort}

\paragraph{Loop invariant.}
At the start of outer-loop iteration $i$, the prefix $A[1 \ldots i-1]$
contains the $i-1$ smallest elements in nondecreasing order, each in its final
position.

\paragraph{Initialization.}
Before the first iteration, $i=1$ and the prefix is empty. It contains the
zero smallest elements in sorted order, so the invariant holds.

\paragraph{Maintenance.}
The algorithm finds the smallest element in $A[i \ldots n]$ and exchanges it
with $A[i]$. This places the next-smallest element in its final position after
the already sorted prefix, so the invariant holds for the next iteration.

\paragraph{Termination.}
After the first $n-1$ iterations, the prefix contains the $n-1$ smallest
elements in sorted order. The one remaining element must be the largest, so the
entire array is sorted. A pass for the $n$th element is unnecessary because no
other element remains to compare or exchange with it.

\paragraph{Running time.}
The algorithm makes the same number of comparisons regardless of input order:
\begin{equation*}
    (n-1) + (n-2) + \cdots + 1 = \frac{n(n-1)}{2}.
\end{equation*}
Therefore, both the best-case and worst-case running times are
$\Theta(n^2)$.
